\chapter{Pika\+Script Binding for LVGL}
\hypertarget{md_managed__components_2lvgl____lvgl_2env__support_2pikascript_2README}{}\label{md_managed__components_2lvgl____lvgl_2env__support_2pikascript_2README}\index{PikaScript Binding for LVGL@{PikaScript Binding for LVGL}}
\label{md_managed__components_2lvgl____lvgl_2env__support_2pikascript_2README_autotoc_md484}%
\Hypertarget{md_managed__components_2lvgl____lvgl_2env__support_2pikascript_2README_autotoc_md484}%
 \href{https://github.com/pikasTech/pikascript}{\texttt{ Pika\+Script}} is an ultralightweight python engine that can run with 4KB of RAM and 32KB of Flash (such as STM32\+G030\+C8 and STM32\+F103\+C8), and is very easy to deploy and expand.

More details to see\+: \href{https://blog.lvgl.io/2022-08-24/pikascript-and-lvgl}{\texttt{ Pika\+Script and lvgl\+: Make Python Lighter, Easier and Smarter}}

The available Python APIs are in the {\ttfamily pika\+\_\+lvgl.\+pyi}, and you need copy the {\ttfamily pika\+\_\+lvgl.\+pyi} to the root path of pikascript, then {\ttfamily import pika\+\_\+lvgl} in {\ttfamily main.\+py}.

The available simulation project on windows\+: \href{https://github.com/pikasTech/lv_pikascript}{\texttt{ https\+://github.\+com/pikas\+Tech/lv\+\_\+pikascript}}

More document about Pika\+Script\+: \href{https://pikadoc.readthedocs.io/en/latest/index.html}{\texttt{ https\+://pikadoc.\+readthedocs.\+io/en/latest/index.\+html}} 